\section{Conclusion}\label{sec:conclusion}

This paper develops an asset-pricing model in which investors search for return predictability in a high-dimensional information environment using a LASSO-based learning rule.
Agents estimate a penalized linear forecasting model, form expectations based on the selected predictors, and these beliefs feed back into prices through a nonlinear asset-pricing equation.
Because perceived and actual laws of motion differ, learning does not converge to rational expectations; instead, the economy settles into a cross-moment consistent equilibrium (CMCE).
The central implication is that the act of searching for predictability can itself generate sparse and short-lived windows of return predictability.
In finite samples, agents occasionally select predictors by chance.
When this occurs, trading on the perceived signal feeds back into prices, making the predictability temporarily self-fulfilling.
As new data arrive and older observations drop out of the agents’ effective memory, coefficient estimates revert back to zero, and the induced predictability disappears.
The theory yields tractable conditions linking the frequency of selection events to fundamentals and features of the information environment.
Empirically, we implement the model on daily returns for 1{,}620 NYSE-listed stocks from 2010--2017 using 180 Wall Street Journal news-attention innovations as predictors.
Consistent with the model, predictability is episodic and sparse, and selection occurs more frequently for higher-volatility stocks.


We plan to extend our analysis along several dimensions.
The current model features a representative agent; a natural next step is to introduce heterogeneity in learning behavior.
On the empirical side, the model cannot explain heterogeneity in the probability that different features are selected, beyond what is implied by their statistical properties. 
The data, however, exhibit systematic cross-feature variation in selection frequencies. 
We therefore test whether this variation reflects an additional, non-statistical component of behavior. 
Specifically, we hypothesize that, after estimating the statistical model, agents are more likely to act on its signals when the resulting predictions are consistent with their pre-existing narratives about how the stock market works and about plausible causal links between specific features and asset prices.



% =====================================================

